\documentclass[11pt]{article}

\usepackage[margin=1in]{geometry}
\usepackage{newtxtext}
\usepackage{microtype}
\usepackage{setspace}
\usepackage[hidelinks]{hyperref}

\setlength{\parindent}{0pt}
\setlength{\parskip}{0.7em}

% The counts below are enforced against the manuscript by
% tests/test_manuscript_tier_visibility.py, which computes them from
% manuscript.tex and fails when a stated value drifts from a computed one.
\newcommand{\papertitle}{Playing time tracks recent exposure in fifteen men's and women's football leagues: the denominator gradient and what it does to per-minute injury rates}
\newcommand{\mainwordcount}{4,484}
\newcommand{\abstractwordcount}{248}
\newcommand{\keywordlist}{injury incidence; exposure; workload; epidemiological methods; women's football; soccer}

\begin{document}

\begin{center}
{\Large\bfseries \papertitle\par}
\vspace{1.5em}
{\large Gustavo Pedro Ricou$^{a,*}$ and Nicholas Mahony$^{b}$\par}
\end{center}

\textbf{Affiliations}

$^{a}$ School of Computer Science and Statistics, Trinity College Dublin,
Dublin 2, Ireland.

$^{b}$ Discipline of Anatomy, School of Medicine, Trinity College Dublin,
Dublin 2, Ireland.

Affiliations are those at which the research was conducted.

\textbf{Corresponding author}

Gustavo Pedro Ricou, School of Computer Science and Statistics, Trinity
College Dublin, Dublin 2, Ireland. Email: \href{mailto:pedrorig@tcd.ie}{pedrorig@tcd.ie}.

\textbf{Author email}

Nicholas Mahony: \href{mailto:njmahony@tcd.ie}{njmahony@tcd.ie}.

\textbf{ORCID}

Gustavo Pedro Ricou: \url{https://orcid.org/0009-0001-4196-7213}.

\textbf{Keywords}

\keywordlist

\textbf{Manuscript details}

Article type: Original Investigation. Main-text word count: \mainwordcount\
(including section headings; excluding abstract, references, tables, figure
captions and declarations). Abstract word count: \abstractwordcount\ (unstructured).
Tables: 3. Figures: 5. Supplementary file: 1. Appendix: 1 (full statistical output).
Preprint: this work is available on arXiv as arXiv:2608.11228
(\url{https://arxiv.org/abs/2608.11228}). The version posted there at the time
of writing (v2) predates the fifteen-league extension: it reports eight men's
leagues and states that the within-starter pattern held without exception. This
manuscript supersedes that claim --- the restriction fails in two of seven
women's leagues --- and a replacement carrying the present analysis has been
submitted as a new version of the same record, so the linked page shows the
superseding version once arXiv announces it.

\textbf{Author contributions (CRediT)}

Gustavo Pedro Ricou: Conceptualization; Data curation; Formal analysis;
Investigation; Methodology; Software; Validation; Visualization; Writing ---
original draft; Writing --- review \& editing. Nicholas Mahony: Methodology;
Writing --- review \& editing.

\textbf{Acknowledgements}

The authors thank James Ng, School of Computer Science and Statistics,
Trinity College Dublin, for statistical feedback on observational framing,
selection bias and internal stability checks. The authors thank Adam Field,
Manchester Metropolitan University, for comments on fixture-congestion
framing, practitioner readability and cautious interpretation. The authors
thank Rilind Ob\"ertinca, University of Innsbruck, for comments on
sports-medicine terminology, presentation and interpretation.

\textbf{Funding}

This research received no specific grant from any funding agency in the
public, commercial, or not-for-profit sectors.

\textbf{Disclosure of interest}

The authors report there are no competing interests to declare.

\textbf{Declaration of generative AI use}

In accordance with the journal's policy on generative artificial intelligence,
Anthropic's Claude was used under author direction to assist with
data-acquisition tooling, analysis code and its verification, literature
searching, and the drafting and editing of manuscript text. The study design,
the scientific claims and their interpretation are the authors'. All reported
quantities are reproducibly computed from the committed data by the archived
code and were verified by the authors, who take full responsibility for the
content of this article.

\textbf{Ethics approval and informed consent}

The study reused exclusively public, pre-existing records, with no
recruitment, contact, intervention or access to private medical records, so no
primary data were collected from human participants and informed consent was
not applicable. No institutional ethics approval number can be provided:
Trinity College Dublin guidance does not permit retrospective ethical approval
of secondary research on already-public data, no determination was sought
before the analysis began, and none can now be issued. The authors will answer
any further enquiry the editorial office wishes to make.

\textbf{Data availability, code availability and data deposition}

The derived, de-identified data and the analysis code that produced every
reported number are deposited together, with the test suite and the hand-built
audit evidence, at \url{https://doi.org/10.5281/zenodo.21940596}, which
resolves to the most recent version of the record. The development repository
is \url{https://github.com/Gustolandia/football-denominator-gradient}. The
complete original output of the statistical software for all 30 reported
gradient models is supplied as Appendix A. Identified source snapshots are not
redistributed, because provider terms and data-protection obligations do not
permit it; their reconstruction is documented in the archived code.

\end{document}
